\documentclass[12pt,letterpaper]{report}

% ---------- Codificación, idioma y márgenes ----------
\usepackage[utf8]{inputenc}
\usepackage[T1]{fontenc}
\usepackage[spanish,es-nodecimaldot,es-noshorthands]{babel}
\usepackage[margin=2.5cm]{geometry}
\usepackage{setspace}
\onehalfspacing

% ---------- Paquetes ----------
\usepackage{graphicx}
\usepackage{float}
\usepackage{array}
\usepackage{booktabs}
\usepackage{longtable}
\usepackage[table]{xcolor}
\usepackage{amsmath,amssymb}
\usepackage{tabularx}
\usepackage{makecell}
\usepackage{seqsplit}
\usepackage{listings}
\usepackage{hyperref}
\hypersetup{hidelinks, colorlinks=false}

% ---------- TikZ para Diagramas UML 2.5 y Grafos de McCabe ----------
\usepackage{tikz}
\usetikzlibrary{arrows.meta,calc,positioning,shapes.geometric,shapes.multipart,fit,backgrounds,babel}

\renewcommand{\thefigure}{\arabic{figure}}
\renewcommand{\thetable}{\arabic{table}}
\AtBeginDocument{\renewcommand{\figurename}{Figura}\renewcommand{\tablename}{Tabla}}
\newcolumntype{L}[1]{>{\raggedright\arraybackslash}p{#1}}
\newcolumntype{C}[1]{>{\centering\arraybackslash}p{#1}}
\newcolumntype{Y}{>{\raggedright\arraybackslash}X}
\newcolumntype{R}[1]{>{\raggedright\arraybackslash}p{\dimexpr\linewidth-(#1)\relax}}

\pagestyle{plain}

% ---------- Configuración de listados de código ----------
\definecolor{codebg}{RGB}{247,248,250}
\definecolor{codeframe}{RGB}{200,205,212}
\definecolor{codekw}{RGB}{0,70,155}
\definecolor{codecmt}{RGB}{70,120,80}
\definecolor{codestr}{RGB}{160,40,40}

\lstdefinestyle{pythonstyle}{
    backgroundcolor=\color{codebg},
    frame=single,
    rulecolor=\color{codeframe},
    basicstyle=\ttfamily\scriptsize\linespread{1.05}\selectfont,
    keywordstyle=\color{codekw}\bfseries,
    commentstyle=\color{codecmt}\itshape,
    stringstyle=\color{codestr},
    numbers=left,
    numberstyle=\tiny\color{gray},
    numbersep=6pt,
    breaklines=true,
    breakatwhitespace=true,
    showstringspaces=false,
    tabsize=4,
    literate={á}{{\'a}}1 {é}{{\'e}}1 {í}{{\'i}}1 {ó}{{\'o}}1 {ú}{{\'u}}1
             {Á}{{\'A}}1 {É}{{\'E}}1 {Í}{{\'I}}1 {Ó}{{\'O}}1 {Ú}{{\'U}}1
             {ñ}{{\~n}}1 {Ñ}{{\~N}}1 {°}{{$^\circ$}}1
}

\begin{document}

% =====================================================================
% PORTADA INSTITUCIONAL
% =====================================================================
\begin{center}
    \parbox[c][\textheight][t]{\textwidth}{
        \begin{center}
            \includegraphics[width=4.5cm]{./figures/sergio_arboleda.png}\\[0.4cm]
            {\normalsize UNIVERSIDAD SERGIO ARBOLEDA}\\ \vspace{-0.1cm}
            {\normalsize ESCUELA DE CIENCIAS EXACTAS E INGENIERÍA}\\ \vspace{-0.1cm}
            {\normalsize CIENCIAS DE LA COMPUTACIÓN E INTELIGENCIA ARTIFICIAL}\\[1.8cm]
            {\large \textbf{PROYECTO FINAL AGRITECH PRECISION — ENTREGA FINAL (SESIÓN 6)}}\\[0.3cm]
            {\normalsize \textbf{DISEÑO ARQUITECTÓNICO UML 2.5, INFORME SQA (ISO/IEC 25010), COMPLEJIDAD DE MCCABE, ÍNDICE DRE Y PROTOTIPO CLOUD}}\\[1.0cm]
            {\normalsize \textbf{JUAN DAVID LÓPEZ YARA}}\\
            {\normalsize \textbf{SANTIAGO JOSÉ RODRÍGUEZ GUTIÉRREZ}}\\
            {\normalsize \textbf{CAMILO ANDRÉS DÍAZ GARCÍA}}\\
            {\normalsize \textbf{ROGER ADRIÁN BOHÓRQUEZ MORENO}}
            \vfill
            {\normalsize \textbf{INGENIERÍA DE SOFTWARE I (PCIA5008)}}\\
            {\normalsize \textbf{JUAN DAVID RODRÍGUEZ MARTÍNEZ}}\\
            BOGOTÁ, NOVIEMBRE 2026
        \end{center}
    }
    \thispagestyle{empty}
\end{center}
\newpage

\pagenumbering{arabic}
\setcounter{page}{1}

% =====================================================================
% SECCIÓN 1: INGENIERÍA DE REQUISITOS (ISO/IEC/IEEE 29148) Y TRAZABILIDAD
% =====================================================================
\section*{1. Ingeniería de Requisitos (ISO/IEC/IEEE 29148) y Modelo de Dominio}

En estricta concordancia con el documento de línea base corregido en la Entrega Parcial (\texttt{DOC-20261112-WA0002}), la arquitectura de software de \textbf{AgriTech Precision} preserva coherencia nominal absoluta entre los requerimientos funcionales y no funcionales, el glosario del dominio, los modelos estructurales/dinámicos en UML 2.5 y la implementación ejecutable.

\subsection*{1.1. Lenguaje Ubicuo y Glosario del Dominio}

La Tabla~\ref{tab:glosario} consolida los conceptos nucleares del dominio junto con sus invariantes de negocio verificables en código.

\begin{table}[H]
    \centering
    \small
    \renewcommand{\arraystretch}{1.3}
    \rowcolors{2}{gray!12}{white}
    \begin{tabularx}{\linewidth}{|L{3.5cm}|L{3.2cm}|Y|}
        \hline
        \rowcolor{gray!40}
        \textbf{Concepto Canónico} & \textbf{Estereotipo / Rol} & \textbf{Definición e Invariante de Dominio} \\
        \hline
        \texttt{SectorAgricola} & Agregado Raíz & Parcela del cultivo identificada por \texttt{idSector}, coordenadas (\texttt{latitud}, \texttt{longitud}) y umbrales de humedad. Compone fuertemente a \texttt{ZonaCultivo}, \texttt{SensorSuelo}, \texttt{CicloRiego}, \texttt{LecturaTelemetria} y \texttt{AlertaTemprana}. Invariante: $0 < \text{umbralCritico} < \text{umbralMinimo} < \text{umbralObjetivo} \le 100\%$ (por defecto: $15\%$, $30\%$ y $60\%$). \\
        \hline
        \texttt{CicloRiego} & Entidad Principal (Máquina de Estados) & Proceso transaccional de irrigación asociado a un \texttt{SectorAgricola}. Registra \texttt{fechaInicio}, \texttt{fechaFin}, \texttt{duracionSegundos}, \texttt{laminaAguaMm}, \texttt{estado} y \texttt{motivoCierre}. Invariante de seguridad: $\text{duracionSegundos} \le 2700\text{ s}$ ($45\text{ min}$). \\
        \hline
        \texttt{SensorSuelo} (\texttt{NodoSensorIoT}) & Entidad de Dominio / Dispositivo IoT & Representación de los nodos sensores (\texttt{NODO-01}, \texttt{NODO-02}, \texttt{NODO-03}) autenticados mediante \texttt{tokenDispositivo} (\texttt{SEG-02}). Solo procesa lecturas cuando $\text{encendido} = \texttt{true}$. \\
        \hline
        \texttt{LecturaTelemetria} & Objeto de Valor Inmutable & Trama enviada por un \texttt{NodoSensorIoT} con humedad y temperatura del suelo. Invariante de validez física: $\text{humedadSuelo} \in [0\%, 100\%]$ y $\text{temperaturaSuelo} \in [-10\,^\circ\text{C}, 60\,^\circ\text{C}]$. \\
        \hline
        \texttt{PronosticoMeteo} & Objeto de Valor Inmutable & Datos climáticos de las próximas $6\text{ horas}$ suministrados por \texttt{APIMeteo}. Vigencia en caché local de $3\text{ horas}$ (\texttt{RNF-03}). \\
        \hline
        \texttt{AlertaTemprana} & Objeto de Valor Inmutable & Notificación por \texttt{RIESGO\_FITOSANITARIO}, \texttt{ESTRES\_TERMICO}, \texttt{RIESGO\_HELADA} o \texttt{FALLO\_VALVULA}. Ventana de supresión de duplicados de $12\text{ horas}$ por sector y tipo (\texttt{RF-06}). \\
        \hline
    \end{tabularx}
    \caption{Glosario canónico e invariantes del dominio AgriTech Precision.}
    \label{tab:glosario}
\end{table}

\subsection*{1.2. Síntesis de Requerimientos Funcionales (RF) y No Funcionales (RNF)}

La Tabla~\ref{tab:rf-sintesis} presenta los cinco Requerimientos Funcionales centrales que gobiernan el flujo automatizado del prototipo (junto con los requerimientos complementarios \texttt{RF-06} a \texttt{RF-09} del SRS), redactados de manera atómica y cuantificable bajo el estándar \textbf{ISO/IEC/IEEE 29148}.

\begin{table}[H]
    \centering
    \small
    \renewcommand{\arraystretch}{1.3}
    \rowcolors{2}{gray!12}{white}
    \begin{tabularx}{\linewidth}{|L{1.3cm}|Y|L{2.3cm}|L{1.6cm}|}
        \hline
        \rowcolor{gray!40}
        \textbf{ID} & \textbf{Especificación Atómica y Cuantitativa (ISO/IEC/IEEE 29148)} & \textbf{Actor} & \textbf{Prioridad} \\
        \hline
        \textbf{RF-01} & El sistema deberá validar cada trama \texttt{\{id\_nodo, humedad, temperatura, timestamp\}} verificando que $\text{humedad} \in [0\%, 100\%]$ y $\text{temperatura} \in [-10\,^\circ\text{C}, 60\,^\circ\text{C}]$, almacenar las lecturas válidas y mostrar por cada \texttt{SectorAgricola} la medición actual, el estado del nodo y de la electroválvula, y el histórico de los últimos $30\text{ días}$. & Gestor de sensores, Técnico agrícola & Must \\
        \hline
        \textbf{RF-02} & El sistema deberá permitir al Gestor de sensores encender y apagar cada \texttt{NodoSensorIoT} (\texttt{NODO-01}, \texttt{02}, \texttt{03}) de forma individual; un nodo con $\text{encendido} = \texttt{false}$ no deberá almacenar lecturas ni disparar evaluaciones de riego. & Gestor de sensores & Must \\
        \hline
        \textbf{RF-03} & El sistema deberá generar y exportar un informe estructurado en formato Excel (\texttt{.xlsx}) para un rango de fechas $[\text{fechaInicio}, \text{fechaFin}]$, conteniendo por fila: \texttt{sector}, \texttt{nodo}, \texttt{fecha\_hora}, \texttt{humedad} y \texttt{temperatura}. & Gestor de sensores & Must \\
        \hline
        \textbf{RF-04} & El sistema deberá ordenar la apertura de la electroválvula (\texttt{activarValvula}) cuando $\text{humedadSuelo} < \text{umbralMinimo}$ ($30\%$ por defecto) y su cierre (\texttt{desactivarValvula}) cuando $\text{humedadSuelo} \ge \text{umbralObjetivo}$ ($60\%$ por defecto) o al cumplirse el temporizador de seguridad $\texttt{after(45 min)}$ ($2700\text{ s}$), registrando inicio, fin, duración y motivo de cierre. & Sistema (automático) & Must \\
        \hline
        \textbf{RF-05} & El sistema deberá consultar a \texttt{APIMeteo} el pronóstico de las próximas $6\text{ horas}$, posponer el riego cuando $\text{probabilidadLluvia} > 70\%$ y la humedad no sea crítica ($\text{humedadSuelo} \ge \text{umbralCritico} = 15\%$), y emitir alertas tempranas (\textbf{RF-06}) ante riesgo fitosanitario ($\text{HR} > 85\%$ y $T \in [18, 28]\,^\circ\text{C}$ por $\ge 6\text{ h}$), estrés térmico ($T_{\text{suelo}} \ge 35\,^\circ\text{C}$ por $\ge 2\text{ h}$) o helada ($T_{\min} \le 2\,^\circ\text{C}$) con supresión de $12\text{ h}$. & Gestor, Técnico agrícola, Sistema & Must \\
        \hline
    \end{tabularx}
    \caption{Requerimientos Funcionales centrales cuantificados bajo ISO/IEC/IEEE 29148.}
    \label{tab:rf-sintesis}
\end{table}

Asimismo, la Tabla~\ref{tab:rnf-sintesis} establece los cinco Requerimientos No Funcionales cuantificables clasificados por característica de calidad según la norma \textbf{ISO/IEC 25010}.

\begin{table}[H]
    \centering
    \small
    \renewcommand{\arraystretch}{1.3}
    \rowcolors{2}{gray!12}{white}
    \begin{tabularx}{\linewidth}{|L{1.4cm}|L{3.1cm}|Y|L{3.2cm}|}
        \hline
        \rowcolor{gray!40}
        \textbf{ID} & \textbf{Característica ISO/IEC 25010} & \textbf{Criterio Cuantitativo de Aceptación} & \textbf{Método de Verificación} \\
        \hline
        \textbf{RNF-01} & Eficiencia de desempeño (Comportamiento temporal y Capacidad) & Procesamiento completo de telemetría (validación, persistencia y evaluación de riego) con latencia $p_{95} \le 2.0\text{ s}$ y emisión de orden a \texttt{ServidorValvulas} en $\le 3.0\text{ s}$, soportando hasta $50$ nodos concurrentes enviando lecturas cada $10\text{ s}$ sin pérdida de paquetes (\textbf{RNF-02}). & Prueba de estrés con $1\,000$ tramas simuladas y medición de percentil 95. \\
        \hline
        \textbf{RNF-03} & Fiabilidad (Disponibilidad y Tolerancia a fallos) & Disponibilidad mensual $\ge 99.0\%$. Timeout hacia \texttt{APIMeteo} de $5.0\text{ s}$ con fallback a caché local vigente ($\le 3\text{ h}$) o decisión autónoma por lectura local. Timeout hacia \texttt{ServidorValvulas} de $3.0\text{ s}$ con hasta $3$ reintentos automáticos (\texttt{IE-03}). & Prueba de inyección de fallos con mocks de \texttt{APIMeteo} y \texttt{ServidorValvulas} caídos. \\
        \hline
        \textbf{RNF-04} & Usabilidad y Portabilidad (Adaptabilidad) & Interfaz web responsiva adaptable a anchos entre $360\text{ px}$ y $1920\text{ px}$ sin desplazamiento horizontal; ejecución de tareas principales en $\le 3\text{ clics}$ y despliegue contenerizado (\texttt{Dockerfile}) portable entre entornos Cloud. & Inspección en 3 resoluciones y construcción de imagen Docker reproducible. \\
        \hline
        \textbf{RNF-05} & Seguridad (Confidencialidad, Integridad y Autenticidad) & Autenticación de nodos mediante \texttt{tokenDispositivo} único revocable (\texttt{SEG-02}, rechazo HTTP \texttt{401}), control de acceso por rol con rechazo HTTP \texttt{403} (\texttt{SEG-01}) y registro del $100\%$ de transacciones críticas en bitácora de auditoría (\textbf{RNF-07}, \texttt{SEG-04}). & Pruebas de penetración funcional sobre endpoints con tokens inválidos y roles sin privilegio. \\
        \hline
        \textbf{RNF-06} & Mantenibilidad (Modificabilidad y Capacidad de prueba) & Desacoplamiento de sistemas externos mediante un único adaptador por puerto (\texttt{IPuertoAPIMeteo}, \texttt{IPuertoServidorValvulas}); incorporación de nuevos nodos en $\le 10\text{ min}$ por configuración; complejidad ciclomática $v(G) \le 10$ en todos los métodos y cobertura de pruebas unitarias $\ge 70\%$. & Auditoría estática de grafos de flujo de control (McCabe) y suite \texttt{unittest}/\texttt{pytest}. \\
        \hline
    \end{tabularx}
    \caption{Requerimientos No Funcionales cuantificables bajo el estándar ISO/IEC 25010.}
    \label{tab:rnf-sintesis}
\end{table}

\subsection*{1.3. Matriz de Trazabilidad (Requisitos $\leftrightarrow$ Diseño $\leftrightarrow$ Pruebas)}

La Tabla~\ref{tab:trazabilidad-diseno} vincula unívocamente cada requerimiento e historia de usuario con los módulos arquitectónicos y los casos de prueba que lo verifican.

\begin{table}[H]
    \centering
    \footnotesize
    \renewcommand{\arraystretch}{1.3}
    \rowcolors{2}{gray!12}{white}
    \begin{tabularx}{\linewidth}{|L{1.6cm}|L{1.8cm}|L{4.2cm}|L{4.2cm}|Y|}
        \hline
        \rowcolor{gray!40}
        \textbf{Req.} & \textbf{HU / CU} & \textbf{Módulo / Clase Responsable} & \textbf{Método / Signatura Exacta} & \textbf{Caso de Prueba Automatizado} \\
        \hline
        RF-01, RF-07 & HU-02, HU-03, CU-03 & \texttt{SensorSuelo}, \texttt{ServicioMonitoreoRiego} & \texttt{validarTrama(humedad, temperatura)}, \texttt{procesarTelemetria(...)} & \texttt{test\_validacion\_rango\_fisico\_sensor\_suelo} \\
        \hline
        RF-02, SEG-01 & HU-01, CU-02 & \texttt{ControladorTelemetriaREST}, \texttt{ServicioMonitoreoRiego} & \texttt{putEstadoNodo(...)}, \texttt{cambiarEstadoNodo(...)} & \texttt{test\_seguridad\_token\_invalido\_401\_y\_rol\_no\_autorizado\_403} \\
        \hline
        RF-03 & HU-08, CU-06 & \texttt{ControladorTelemetriaREST} & \texttt{exportarFilasInformeExcel(sectores, fechaInicio, fechaFin)} & Verificación de estructura tabular de exportación \\
        \hline
        RF-04, RF-08 & HU-04, HU-09, HU-11, CU-04 & \texttt{CicloRiego}, \texttt{EstrategiaRiegoAgronomico}, \texttt{AdaptadorServidorValvulas} & \texttt{calcularLaminaRiego(...)}, \texttt{activarValvula(...)}, \texttt{verificarTimeoutRiego(...)} & \texttt{test\_camino\_1} a \texttt{test\_camino\_6}, \texttt{test\_cierre\_seguridad\_temporizador\_after\_45\_min} \\
        \hline
        RF-05, RF-06, RNF-03 & HU-05, HU-06, HU-07, CU-05, CU-07 & \texttt{ServicioMonitoreoRiego}, \texttt{AdaptadorAPIMeteo} & \texttt{consultarPronostico(...)}, \texttt{evaluarRiesgoFitosanitario(...)} & \texttt{test\_mccabe\_9\_caminos\_evaluar\_riesgo\_fitosanitario}, \texttt{test\_resiliencia\_cache\_apimeteo} \\
        \hline
    \end{tabularx}
    \caption{Matriz de trazabilidad extremo a extremo entre SRS, Arquitectura y Pruebas.}
    \label{tab:trazabilidad-diseno}
\end{table}

% =====================================================================
% SECCIÓN 2: MODELADO DE INTERACCIÓN UML 2.5 (DIAGRAMA DE SECUENCIA)
% =====================================================================
\newpage
\section*{2. Modelado de Interacción UML 2.5 — Diagrama de Secuencia}

La Figura~\ref{fig:secuencia} modela la interacción temporal completa del flujo principal (\texttt{CU-04 Activar riego automático} e integración con evaluación fitosanitaria) siguiendo estrictamente la sintaxis de \textbf{OMG UML 2.5}:
\begin{itemize}
    \item \textbf{Líneas de vida estereotipadas:} \texttt{<<external system>> :NodoSensorIoT}, \texttt{<<boundary>> :ControladorTelemetriaREST}, \texttt{<<control>> :ServicioMonitoreoRiego}, la entidad \texttt{:SectorAgricola / :CicloRiego}, \texttt{<<external system>> :AdaptadorAPIMeteo} y \texttt{<<external system>> :AdaptadorServidorValvulas}.
    \item \textbf{Mensajes síncronos vs. asíncronos:} Los mensajes asíncronos se representan con punta de flecha abierta (\texttt{->>}, envío de telemetría por red y despacho de alertas/auditoría), mientras que las invocaciones bloqueantes entre capas usan punta de flecha cerrada rellena (\texttt{->}) con su respectivo mensaje de retorno en línea discontinua (\texttt{-->}).
    \item \textbf{Fragmentos combinados anidados:}
    \begin{enumerate}
        \item \texttt{loop [cada intervalo\_s = 10 s mientras sensor.encendido == true]}: controla el ciclo periódico de ingesta de telemetría.
        \item \texttt{par [Evaluación concurrente]}: ejecuta en paralelo la consulta de pronóstico meteorológico + evaluación de riesgo fitosanitario (\texttt{evaluarRiesgoFitosanitario}) y la evaluación de necesidad de riego (\texttt{calcularLaminaRiego}).
        \item \texttt{alt [humedadSuelo < umbralMinimo]}: bifurca la decisión de riego entre posponer por lluvia (\texttt{[probLluvia > 70\% y humedadSuelo >= 15\%]}), activar la electroválvula (\texttt{[laminaMm > 0]}) con vigilancia del temporizador \texttt{after(45 min)}, o mantener el estado inactivo (\texttt{[humedadSuelo >= umbralMinimo]}).
    \end{enumerate}
\end{itemize}

\begin{figure}[H]
    \centering
    {\linespread{1}\selectfont
    \resizebox{\linewidth}{!}{%
    \begin{tikzpicture}[
        lifeline/.style={draw=blue!55!black, thick, fill=blue!8, rectangle, rounded corners=3pt,
                         minimum width=2.65cm, minimum height=1.15cm, align=center, font=\scriptsize},
        extlife/.style={draw=orange!65!black, thick, fill=orange!10, rectangle, rounded corners=3pt,
                        minimum width=2.65cm, minimum height=1.15cm, align=center, font=\scriptsize},
        syncmsg/.style={-{Stealth[length=2.4mm, width=1.8mm]}, thick, draw=black},
        asyncmsg/.style={-{Straight Barb[length=2.5mm, width=2.2mm]}, thick, draw=blue!65!black},
        retmsg/.style={-{Straight Barb[length=2.2mm]}, dashed, semithick, draw=black!75},
        actbox/.style={draw=black!80, fill=gray!20, rectangle, minimum width=0.22cm},
        frag/.style={draw=black!75, thick, rectangle},
        fragtab/.style={draw=black!75, thick, fill=gray!25, font=\scriptsize\bfseries, inner sep=2.5pt, anchor=north west}
    ]
        % Coordenadas X de las líneas de vida
        \def\xNodo{0}
        \def\xCtrl{3.4}
        \def\xSvc{7.0}
        \def\xDom{10.6}
        \def\xMeteo{14.2}
        \def\xValv{17.8}

        % Cabeceras de líneas de vida
        \node[extlife] (nodo) at (\xNodo, 0) {\texttt{<<external system>>}\\\textbf{:NodoSensorIoT}};
        \node[lifeline] (ctrl) at (\xCtrl, 0) {\texttt{<<boundary>>}\\\textbf{:Controlador}\\\textbf{TelemetriaREST}};
        \node[lifeline] (svc) at (\xSvc, 0) {\texttt{<<control>>}\\\textbf{:Servicio}\\\textbf{MonitoreoRiego}};
        \node[lifeline] (dom) at (\xDom, 0) {\textbf{:SectorAgricola}\\{\tiny\textbf{/ :CicloRiego}}};
        \node[extlife] (meteo) at (\xMeteo, 0) {\texttt{<<external system>>}\\\textbf{:Adaptador}\\\textbf{APIMeteo}};
        \node[extlife] (valv) at (\xValv, 0) {\texttt{<<external system>>}\\\textbf{:Adaptador}\\\textbf{ServidorValvulas}};

        % Líneas punteadas verticales
        \foreach \x in {\xNodo, \xCtrl, \xSvc, \xDom, \xMeteo, \xValv} {
            \draw[dashed, gray!70, thick] (\x, -0.6) -- (\x, -18.8);
        }

        % Cajas de activación
        \draw[actbox] (\xCtrl-0.11, -1.6) rectangle (\xCtrl+0.11, -18.0);
        \draw[actbox] (\xSvc-0.11, -2.1) rectangle (\xSvc+0.11, -17.6);
        \draw[actbox] (\xDom-0.11, -2.7) rectangle (\xDom+0.11, -4.0);
        \draw[actbox] (\xMeteo-0.11, -5.2) rectangle (\xMeteo+0.11, -6.3);
        \draw[actbox] (\xDom-0.11, -7.0) rectangle (\xDom+0.11, -9.2);
        \draw[actbox] (\xValv-0.11, -12.6) rectangle (\xValv+0.11, -13.7);
        \draw[actbox] (\xValv-0.11, -15.6) rectangle (\xValv+0.11, -16.5);

        % --- FRAGMENTO COMBINADO LOOP ---
        \draw[frag] (-1.2, -1.0) rectangle (19.1, -18.4);
        \node[fragtab] at (-1.2, -1.0) {loop [cada intervalo\_s = 10 s mientras sensor.encendido == true]};

        % 1. Mensaje asíncrono de telemetría
        \draw[asyncmsg] (\xNodo, -1.6) -- node[above, font=\scriptsize] {1: \texttt{postLecturaTelemetria(payload, token)} {\tiny[async]}} (\xCtrl-0.11, -1.6);

        % 2. Invocación síncrona al servicio
        \draw[syncmsg] (\xCtrl+0.11, -2.1) -- node[above, font=\scriptsize] {2: \texttt{procesarTelemetria(idNodo, hum, temp, token, ts)}} (\xSvc-0.11, -2.1);

        % 3. Validación y registro de lectura
        \draw[syncmsg] (\xSvc+0.11, -2.8) -- node[above, font=\scriptsize] {3: \texttt{validarTrama(hum, temp)} y \texttt{registrarLectura(lectura)}} (\xDom-0.11, -2.8);
        \draw[retmsg] (\xDom-0.11, -3.7) -- node[above, font=\scriptsize] {3.1: \texttt{esValida = true}} (\xSvc+0.11, -3.7);

        % --- FRAGMENTO COMBINADO PAR ---
        \draw[frag, draw=teal!75!black] (5.5, -4.4) rectangle (19.0, -9.5);
        \node[fragtab, fill=teal!20] at (5.5, -4.4) {par [Consulta meteorológica y evaluación fitosanitaria concurrente]};

        % Rama 1 de PAR: Consulta a APIMeteo
        \draw[syncmsg] (\xSvc+0.11, -5.3) -- node[above, font=\scriptsize] {4: \texttt{consultarPronostico(latitud, longitud)}} (\xMeteo-0.11, -5.3);
        \draw[retmsg] (\xMeteo-0.11, -6.2) -- node[above, font=\scriptsize] {4.1: \texttt{PronosticoMeteo(temp, HR, probLluvia, Tmin)}} (\xSvc+0.11, -6.2);

        \draw[dashed, thick, teal!75!black] (5.5, -6.6) -- (19.0, -6.6);
        \node[anchor=north west, font=\scriptsize\bfseries, text=teal!70!black] at (5.6, -6.6) {[Rama concurrente: Evaluación de Riesgo Fitosanitario (RF-06) y Cálculo de Lámina]};

        \draw[syncmsg] (\xSvc+0.11, -7.4) -- node[above, font=\scriptsize] {5: \texttt{evaluarRiesgoFitosanitario(sector, lectura, pronostico)}} (\xDom-0.11, -7.4);
        \draw[asyncmsg] (\xDom-0.11, -8.2) -- node[above, font=\scriptsize] {5.1: \texttt{actualizarAlerta(alerta)} {\tiny[async Observer]}} (\xSvc+0.11, -8.2);
        \draw[syncmsg] (\xSvc+0.11, -8.9) -- node[above, font=\scriptsize] {6: \texttt{calcularLaminaRiego(hum, min, crit, obj, probLluvia, temp)}} (\xDom-0.11, -8.9);

        % --- FRAGMENTO COMBINADO ALT ---
        \draw[frag, draw=purple!75!black] (5.5, -9.9) rectangle (19.0, -17.2);
        \node[fragtab, fill=purple!18] at (5.5, -9.9) {alt [humedadSuelo < umbralMinimo (30\%)]};

        % Sub-rama A: Lluvia prevista y humedad no crítica
        \node[anchor=north west, font=\scriptsize\bfseries, text=purple!80!black] at (5.7, -10.5) {[probLluvia > 70\% AND humedadSuelo >= umbralCritico (15\%)]};
        \draw[syncmsg] (\xSvc+0.11, -11.3) -- node[above, font=\scriptsize] {7a: \texttt{transicionarA(POSPUESTO\_POR\_LLUVIA)}} (\xDom-0.11, -11.3);

        \draw[dashed, thick, purple!75!black] (5.5, -11.8) -- (19.0, -11.8);
        \node[anchor=north west, font=\scriptsize\bfseries, text=purple!80!black] at (5.7, -11.85) {[else: probLluvia <= 70\% OR humedadSuelo < 15\% (Riego Procede)]};

        % Apertura de electroválvula
        \draw[syncmsg] (\xSvc+0.11, -12.7) -- node[above, font=\scriptsize] {7b: \texttt{activarValvula(sectorId, duracionSegundos=2700)}} (\xValv-0.11, -12.7);
        \draw[retmsg] (\xValv-0.11, -13.5) -- node[above, font=\scriptsize] {7b.1: \texttt{EstadoValvulaDTO(estado="abierta")}} (\xSvc+0.11, -13.5);
        \draw[syncmsg] (\xSvc+0.11, -14.2) -- node[above, font=\scriptsize] {8: \texttt{transicionarA(IRRIGANDO\_FLUJO\_ACTIVO)}} (\xDom-0.11, -14.2);

        % Cierre por objetivo o temporizador after(45 min)
        \draw[syncmsg] (\xSvc+0.11, -15.6) -- node[above, font=\scriptsize] {9: \texttt{desactivarValvula(sectorId)} {\tiny[humedad >= 60\% OR after(45 min)]}} (\xValv-0.11, -15.6);
        \draw[retmsg] (\xValv-0.11, -16.3) -- node[above, font=\scriptsize] {9.1: \texttt{EstadoValvulaDTO(estado="cerrada")}} (\xSvc+0.11, -16.3);

        \draw[dashed, thick, purple!75!black] (5.5, -16.6) -- (19.0, -16.6);
        \node[anchor=north west, font=\scriptsize\bfseries, text=purple!80!black] at (5.7, -16.62) {[else: humedadSuelo >= umbralMinimo (30\%)] -> mantener INACTIVO};

        % Retorno final al controlador
        \draw[retmsg] (\xSvc-0.11, -17.5) -- node[above, font=\scriptsize] {10: \texttt{ResultadoProcesamiento(codigoHttp=201, estadoValvula, laminaMm)}} (\xCtrl+0.11, -17.5);
        \draw[retmsg] (\xCtrl-0.11, -17.9) -- node[above, font=\scriptsize] {11: \texttt{HTTP 201 Created \{estado\_sector, estado\_valvula\}}} (\xNodo, -17.9);
    \end{tikzpicture}}}
    \small
    \caption{Diagrama de Secuencia UML 2.5 con estereotipos \texttt{<<boundary>>}, \texttt{<<control>>}, \texttt{<<external system>>} y fragmentos combinados \texttt{loop}, \texttt{par} y \texttt{alt}.}
    \label{fig:secuencia}
\end{figure}

% =====================================================================
% SECCIÓN 3: DISEÑO ESTRUCTURAL Y MÁQUINA DE ESTADOS UML 2.5
% =====================================================================
\newpage
\section*{3. Diseño Estructural y Máquina de Estados UML 2.5}

\subsection*{3.1. Principios de Diseño: Módulos Profundos, Patrones GoF y Composición vs. Agregación}

El diseño estructural aplica la disciplina de \textbf{módulos profundos} (\textit{Deep Modules}) y tres patrones de diseño GoF:
\begin{enumerate}
    \item \textbf{Módulo Profundo (\texttt{ServicioMonitoreoRiego}):} Concentra toda la orquestación del dominio detrás de una interfaz pública reducida (5 métodos públicos en su \textit{Seam} de aplicación), ocultando a los llamadores la validación de tramas, la gestión de caché meteorológico, los reintentos de red y las transiciones de estado.
    \item \textbf{Patrones GoF Aplicados:}
    \begin{itemize}
        \item \textbf{Strategy (\texttt{IEstrategiaRiego} $\leftarrow$ \texttt{EstrategiaRiegoAgronomico}):} Encapsula la fórmula de cálculo de lámina de riego (\texttt{calcularLaminaRiego}), permitiendo variar el modelo agronómico sin alterar \texttt{ServicioMonitoreoRiego}.
        \item \textbf{Observer (\texttt{IObservadorAlerta} $\leftarrow$ \texttt{BitacoraAuditoriaObserver}):} Desacopla la detección de alertas fitosanitarias y eventos de seguridad de sus canales de notificación y registro de auditoría (\texttt{RNF-07}).
        \item \textbf{Adapter (\texttt{IPuertoAPIMeteo} $\leftarrow$ \texttt{AdaptadorAPIMeteo}, \texttt{IPuertoServidorValvulas} $\leftarrow$ \texttt{AdaptadorServidorValvulas}):} Aísla cada dependencia externa tras un único módulo adaptador (\texttt{RNF-06}).
    \end{itemize}
    \item \textbf{Distinción Estricta entre Composición Fuerte y Agregación:}
    \begin{itemize}
        \item \textbf{Composición Fuerte (\texttt{*--}, rombo relleno negro):} \texttt{SectorAgricola} es propietario exclusivo del ciclo de vida de sus instancias de \texttt{ZonaCultivo} ($1..*$), \texttt{SensorSuelo} ($1..*$), \texttt{CicloRiego} ($0..*$), \texttt{LecturaTelemetria} ($0..*$) y \texttt{AlertaTemprana} ($0..*$). Si un \texttt{SectorAgricola} se elimina, sus partes dejan de existir.
        \item \textbf{Agregación (\texttt{o--}, rombo blanco vacío):} \texttt{ServicioMonitoreoRiego} mantiene referencias por agregación a la colección de suscriptores \texttt{IObservadorAlerta} ($0..*$) y a la estrategia \texttt{IEstrategiaRiego} ($1$), los cuales tienen ciclo de vida independiente y pueden compartirse entre servicios.
    \end{itemize}
\end{enumerate}

\begin{figure}[H]
    \centering
    {\linespread{0.95}\selectfont
    \resizebox{\linewidth}{!}{%
    \begin{tikzpicture}[
        classbox/.style={draw=blue!55!black, thick, fill=white, rectangle, rounded corners=2pt,
                         font=\scriptsize, align=left, inner sep=4pt},
        ifacebox/.style={draw=teal!60!black, thick, fill=teal!5, rectangle, rounded corners=2pt,
                         font=\scriptsize, align=left, inner sep=4pt},
        comp/.style={{Diamond[length=3.2mm, width=2.2mm]}-{Straight Barb[length=2.2mm]}, thick},
        aggr/.style={{Diamond[open, length=3.2mm, width=2.2mm]}-{Straight Barb[length=2.2mm]}, thick},
        impl/.style={-{Triangle[open, length=2.8mm, width=2.4mm]}, dashed, thick},
        assoc/.style={-{Straight Barb[length=2.2mm]}, thick}
    ]
        % Controlador Boundary
        \node[classbox] (ctrl) at (0, 6.2) {
            \texttt{<<boundary>>}\\\textbf{ControladorTelemetriaREST}\\[-2pt]\rule{5.2cm}{0.35pt}\\
            \texttt{- \_servicio: ServicioMonitoreoRiego}\\[-2pt]\rule{5.2cm}{0.35pt}\\
            \texttt{+ postLecturaTelemetria(payload: Dict, tokenHeader: str): Dict}\\
            \texttt{+ putEstadoNodo(idNodo: str, encendido: bool, rolHeader: str, email: str): Dict}\\
            \texttt{+ exportarFilasInformeExcel(sectores: List, ini: datetime, fin: datetime): List}
        };

        % Servicio Control (Deep Module)
        \node[classbox] (svc) at (0, 2.2) {
            \texttt{<<control>>}\\\textbf{ServicioMonitoreoRiego}\\[-2pt]\rule{5.2cm}{0.35pt}\\
            \texttt{- \_repositorio: IRepositorioSectores}\\
            \texttt{- \_puertoMeteo: IPuertoAPIMeteo}\\
            \texttt{- \_puertoValvulas: IPuertoServidorValvulas}\\
            \texttt{- \_estrategiaRiego: IEstrategiaRiego}\\
            \texttt{- \_observadores: List[IObservadorAlerta]}\\[-2pt]\rule{5.2cm}{0.35pt}\\
            \texttt{+ procesarTelemetria(idNodo: str, hum: float, temp: float, token: str, ts: datetime): ResultadoProcesamiento}\\
            \texttt{+ evaluarRiesgoFitosanitario(sector: SectorAgricola, lec: LecturaTelemetria, pron: PronosticoMeteo, hFungo: int, hCalor: int): Optional[AlertaTemprana]}\\
            \texttt{+ calcularLaminaRiego(hum: float, min: float, crit: float, obj: float, lluvia: float, temp: float): float}\\
            \texttt{+ verificarTimeoutRiego(idSector: str, segs: int, humActual: float, ts: datetime): Optional[CicloRiego]}\\
            \texttt{+ cambiarEstadoNodo(idNodo: str, encendido: bool, rol: RolUsuario, email: str, ts: datetime): SensorSuelo}
        };

        % Interfaces Strategy, Observer y Ports (Derecha arriba/medio)
        \node[ifacebox] (istrat) at (10.8, 6.3) {
            \texttt{<<interface>>}\\\textbf{IEstrategiaRiego}\\[-2pt]\rule{5.2cm}{0.35pt}\\
            \texttt{+ calcularLaminaRiego(hum: float, min: float,}\\
            \texttt{  crit: float, obj: float, lluvia: float, temp: float): float}
        };

        \node[classbox] (strat) at (10.8, 3.8) {
            \textbf{EstrategiaRiegoAgronomico}\\[-2pt]\rule{5.2cm}{0.35pt}\\
            \texttt{+ calcularLaminaRiego(hum: float, min: float,}\\
            \texttt{  crit: float, obj: float, lluvia: float, temp: float): float}
        };

        \node[ifacebox] (iobs) at (10.8, 1.5) {
            \texttt{<<interface>>}\\\textbf{IObservadorAlerta}\\[-2pt]\rule{5.2cm}{0.35pt}\\
            \texttt{+ actualizarAlerta(alerta: AlertaTemprana): None}\\
            \texttt{+ registrarEventoAuditoria(usr: str, op: str, det: str, ts: datetime): None}
        };

        \node[classbox] (obs) at (10.8, -0.8) {
            \textbf{BitacoraAuditoriaObserver}\\[-2pt]\rule{5.2cm}{0.35pt}\\
            \texttt{+ alertasRecibidas: List[AlertaTemprana]}\\
            \texttt{+ registrosBitacora: List[Dict]}\\[-2pt]\rule{5.2cm}{0.35pt}\\
            \texttt{+ actualizarAlerta(alerta: AlertaTemprana): None}\\
            \texttt{+ registrarEventoAuditoria(...): None}
        };

        % Puertos y Adaptadores Externos (Izquierda medio/abajo)
        \node[ifacebox] (imeteo) at (-7.2, 2.5) {
            \texttt{<<interface>>}\\\textbf{IPuertoAPIMeteo}\\[-2pt]\rule{5.2cm}{0.35pt}\\
            \texttt{+ consultarPronostico(lat: float, lon: float): PronosticoMeteo}
        };

        \node[classbox] (ameteo) at (-7.2, -0.2) {
            \texttt{<<external adapter>>}\\\textbf{AdaptadorAPIMeteo}\\[-2pt]\rule{5.2cm}{0.35pt}\\
            \texttt{+ VIGENCIA\_CACHE\_HORAS: int = 3}\\
            \texttt{+ TIMEOUT\_SEGUNDOS: float = 5.0}\\
            \texttt{- \_cache: Dict[Tuple, PronosticoMeteo]}\\[-2pt]\rule{5.2cm}{0.35pt}\\
            \texttt{+ consultarPronostico(lat: float, lon: float): PronosticoMeteo}
        };

        \node[ifacebox] (ivalv) at (-7.2, -2.8) {
            \texttt{<<interface>>}\\\textbf{IPuertoServidorValvulas}\\[-2pt]\rule{5.2cm}{0.35pt}\\
            \texttt{+ activarValvula(sectorId: str, duracionSegundos: int): EstadoValvulaDTO}\\
            \texttt{+ desactivarValvula(sectorId: str): EstadoValvulaDTO}
        };

        \node[classbox] (avalv) at (-7.2, -5.5) {
            \texttt{<<external adapter>>}\\\textbf{AdaptadorServidorValvulas}\\[-2pt]\rule{5.2cm}{0.35pt}\\
            \texttt{+ MAX\_REINTENTOS: int = 3}\\
            \texttt{+ TIMEOUT\_SEGUNDOS: float = 3.0}\\[-2pt]\rule{5.2cm}{0.35pt}\\
            \texttt{+ activarValvula(sectorId: str, duracionSegundos: int): EstadoValvulaDTO}\\
            \texttt{+ desactivarValvula(sectorId: str): EstadoValvulaDTO}
        };

        % Agregado Raíz SectorAgricola (Centro abajo)
        \node[classbox] (sector) at (0, -3.5) {
            \texttt{<<aggregate root>>}\\\textbf{SectorAgricola}\\[-2pt]\rule{5.2cm}{0.35pt}\\
            \texttt{+ idSector: str}\\
            \texttt{+ nombre: str}\\
            \texttt{+ idValvula: str}\\
            \texttt{+ latitud: float; + longitud: float}\\
            \texttt{+ umbralCritico: float = 15.0}\\
            \texttt{+ umbralMinimo: float = 30.0}\\
            \texttt{+ umbralObjetivo: float = 60.0}\\
            \texttt{+ estadoSector: EstadoSector}\\[-2pt]\rule{5.2cm}{0.35pt}\\
            \texttt{+ configurarUmbrales(critico: float, minimo: float, objetivo: float): None}\\
            \texttt{+ registrarLectura(lectura: LecturaTelemetria): None}\\
            \texttt{+ iniciarCicloRiego(ciclo: CicloRiego): None}\\
            \texttt{+ cerrarCicloActivo(humFinal: float, dur: int, mot: MotivoCierreRiego, ts: datetime): CicloRiego}\\
            \texttt{+ tieneAlertaReciente(tipo: TipoAlerta, actual: datetime, ventanaHoras: int): bool}
        };

        % Entidades compuestas fuertemente por SectorAgricola (Abajo)
        \node[classbox] (zona) at (-6.8, -8.8) {
            \textbf{ZonaCultivo}\\[-2pt]\rule{5.2cm}{0.35pt}\\
            \texttt{+ idZona: str}\\
            \texttt{+ tipoCultivo: str}\\
            \texttt{+ areaMetrosCuadrados: float}\\
            \texttt{+ coeficienteCultivoKc: float}\\[-2pt]\rule{5.2cm}{0.35pt}\\
            \texttt{+ calcularDemandaHidricaBase(): float}
        };

        \node[classbox] (sensor) at (-1.5, -8.8) {
            \textbf{SensorSuelo}\\[-2pt]\rule{5.2cm}{0.35pt}\\
            \texttt{+ idNodo: str}\\
            \texttt{+ idSector: str}\\
            \texttt{+ tokenDispositivo: str}\\
            \texttt{+ encendido: bool = True}\\
            \texttt{+ intervaloSegundos: int = 10}\\[-2pt]\rule{5.2cm}{0.35pt}\\
            \texttt{+ cambiarEstadoOperativo(enc: bool): None}\\
            \texttt{+ validarTrama(hum: float, temp: float): bool}
        };

        \node[classbox] (ciclo) at (4.5, -8.8) {
            \textbf{CicloRiego}\\[-2pt]\rule{5.2cm}{0.35pt}\\
            \texttt{+ idCiclo: str; + idSector: str}\\
            \texttt{+ estado: EstadoCicloRiego}\\
            \texttt{+ fechaInicio: datetime; + fechaFin: Optional[datetime]}\\
            \texttt{+ humedadInicial: float; + humedadFinal: Optional[float]}\\
            \texttt{+ laminaAguaMm: float; + duracionSegundos: int}\\
            \texttt{+ motivoCierre: Optional[MotivoCierreRiego]}\\
            \texttt{+ MAX\_DURACION\_SEGUNDOS: int = 2700}\\[-2pt]\rule{5.2cm}{0.35pt}\\
            \texttt{+ transicionarA(nuevo: EstadoCicloRiego, mot: MotivoCierreRiego): None}\\
            \texttt{+ estaActivo(): bool}\\
            \texttt{+ finalizar(humFin: float, dur: int, mot: MotivoCierreRiego, ts: datetime): None}
        };

        \node[classbox] (alerta) at (10.8, -4.5) {
            \textbf{AlertaTemprana}\\[-2pt]\rule{5.2cm}{0.35pt}\\
            \texttt{+ idAlerta: str}\\
            \texttt{+ idSector: str}\\
            \texttt{+ tipoAlerta: TipoAlerta}\\
            \texttt{+ severidad: str}\\
            \texttt{+ mensaje: str}\\
            \texttt{+ timestamp: datetime}
        };

        % Relaciones
        \draw[assoc] (ctrl.south) -- node[right, font=\scriptsize] {\texttt{1}} (svc.north);
        \draw[aggr] (svc.east) -- node[above, font=\scriptsize, pos=0.5] {agrega \texttt{1}} (istrat.west);
        \draw[impl] (strat.north) -- (istrat.south);
        \draw[aggr] (svc.east) -- node[below, font=\scriptsize, pos=0.5] {agrega \texttt{0..*}} (iobs.west);
        \draw[impl] (obs.north) -- (iobs.south);

        \draw[assoc] (svc.west) -- node[above, font=\scriptsize] {\texttt{1}} (imeteo.east);
        \draw[impl] (ameteo.north) -- (imeteo.south);
        \draw[assoc] (svc.west) -- node[below, font=\scriptsize, pos=0.6] {\texttt{1}} (ivalv.east);
        \draw[impl] (avalv.north) -- (ivalv.south);

        \draw[assoc] (svc.south) -- node[right, font=\scriptsize] {gestiona \texttt{1..*}} (sector.north);

        % Composiciones fuertes (*--)
        \draw[comp] (sector.south west) -- node[left, font=\scriptsize, pos=0.7] {\texttt{1..*} (Composición)} (zona.north);
        \draw[comp] (sector.south) -- node[left, font=\scriptsize, pos=0.7] {\texttt{1..*} (Composición)} (sensor.north);
        \draw[comp] (sector.south east) -- node[right, font=\scriptsize, pos=0.7] {\texttt{0..*} (Composición)} (ciclo.north);
        \draw[comp] (sector.east) -- node[above, font=\scriptsize, pos=0.6] {\texttt{0..*} (Composición)} (alerta.west);
    \end{tikzpicture}}}
    \small
    \caption{Diagrama de Clases de Diseño UML 2.5 con visibilidad, signaturas completas, patrones GoF (\textit{Strategy}, \textit{Observer}, \textit{Adapter}) y distinción explícita entre Composición Fuerte (rombo relleno) y Agregación (rombo vacío).}
    \label{fig:clases}
\end{figure}

\subsection*{3.2. Diagrama de Máquina de Estados UML 2.5 (\texttt{CicloRiego})}

La Figura~\ref{fig:estados} modela el ciclo de vida completo del objeto de negocio principal \texttt{CicloRiego} utilizando la sintaxis formal \textbf{\texttt{Evento [Guarda] / Acción}}, el estado compuesto \textbf{\texttt{REGANDO}} con tres subestados anidados y los temporizadores explícitos \textbf{\texttt{after(3 s)}} y \textbf{\texttt{after(45 min)}}.

\begin{figure}[H]
    \centering
    {\linespread{1}\selectfont
    \resizebox{\linewidth}{!}{%
    \begin{tikzpicture}[
        state/.style={draw=blue!60!black, very thick, fill=blue!7, rectangle, rounded corners=6pt,
                      minimum width=3.4cm, minimum height=1.1cm, align=center, font=\scriptsize},
        substate/.style={draw=teal!65!black, thick, fill=white, rectangle, rounded corners=5pt,
                         minimum width=3.2cm, minimum height=0.95cm, align=center, font=\scriptsize},
        errstate/.style={draw=red!70!black, very thick, fill=red!8, rectangle, rounded corners=6pt,
                         minimum width=3.3cm, minimum height=1.1cm, align=center, font=\scriptsize},
        trans/.style={-{Stealth[length=2.5mm]}, thick, draw=black!85},
        lbl/.style={font=\scriptsize, align=center, inner sep=2pt}
    ]
        % Nodo inicial
        \filldraw[black] (-6.5, 4.5) circle (0.18cm) node (init) {};

        % Estados superiores
        \node[state] (inactivo) at (-3.2, 4.5) {\textbf{INACTIVO}\\[2pt]{\tiny entry / estadoValvula = "cerrada"}};
        \node[state] (eval) at (3.2, 4.5) {\textbf{EVALUANDO\_UMBRAL}\\[2pt]{\tiny do / comparar humedadSuelo vs umbralMinimo}};
        \node[state] (meteo) at (10.2, 4.5) {\textbf{CONSULTANDO\_METEO}\\[2pt]{\tiny do / consultarPronostico(lat, lon)}};
        \node[state] (pospuesto) at (10.2, 1.2) {\textbf{POSPUESTO\_POR\_LLUVIA}\\[2pt]{\tiny entry / registrarMotivo(POSPUESTO\_POR\_LLUVIA)}};

        % Estado compuesto REGANDO con subestados anidados
        \draw[draw=teal!70!black, very thick, rounded corners=9pt, fill=teal!6] (-5.8, -0.3) rectangle (6.8, -5.8);
        \node[font=\footnotesize\bfseries, text=teal!60!black, anchor=north west] at (-5.6, -0.4) {Estado Compuesto: REGANDO (Válvula bajo control activo)};

        \node[substate] (abriendo) at (-3.5, -2.0) {\textbf{ABRIENDO\_VALVULA}\\[2pt]{\tiny entry / activarValvula(id, 2700)}};
        \node[substate] (irrigando) at (3.6, -2.0) {\textbf{IRRIGANDO\_FLUJO\_ACTIVO}\\[2pt]{\tiny do / recibir telemetría periódica}};
        \node[substate] (verif) at (3.6, -4.5) {\textbf{VERIFICANDO\_HUMEDAD}\\\textbf{\_OBJETIVO}};

        % Estados de cierre y terminales
        \node[state] (cerrando) at (-3.5, -7.8) {\textbf{CERRANDO\_VALVULA}\\[2pt]{\tiny entry / desactivarValvula(id)}};
        \node[state] (completado) at (-3.5, -10.5) {\textbf{COMPLETADO}\\[2pt]{\tiny entry / registrarFin(OBJETIVO\_ALCANZADO)}};
        \node[errstate] (timeout) at (3.6, -10.5) {\textbf{ABORTADO\_TIMEOUT}\\[2pt]{\tiny entry / emitirAlerta(TIEMPO\_MAXIMO)}};
        \node[errstate] (fallo) at (10.2, -5.0) {\textbf{FALLO\_VALVULA}\\[2pt]{\tiny entry / emitirAlerta(FALLO\_VALVULA)}};

        % Nodo final
        \draw[thick] (3.6, -12.6) circle (0.24cm) node (final) {};
        \filldraw[black] (3.6, -12.6) circle (0.14cm);

        % Transiciones
        \draw[trans] (init) -- (inactivo);
        \draw[trans] (inactivo.east) -- node[lbl, above] {\texttt{lecturaRecibida [esValida == true]}\\\texttt{/ almacenarLectura()}} (eval.west);
        \draw[trans] (eval.south) to[bend left=18] node[lbl, below] {\texttt{[humedadSuelo >= umbralMinimo (30\%)] / mantenerInactivo()}} (inactivo.south);
        \draw[trans] (eval.east) -- node[lbl, above] {\texttt{[humedadSuelo < umbralMinimo]}\\\texttt{/ iniciarEvaluacionMeteo()}} (meteo.west);

        \draw[trans] (meteo.south) -- node[lbl, right] {\texttt{pronosticoRecibido}\\\texttt{[probLluvia > 70\% AND humedadSuelo >= 15\%]}\\\texttt{/ posponerCiclo()}} (pospuesto.north);

        \draw[trans] (meteo.south west) -- node[lbl, above, sloped] {\texttt{[probLluvia <= 70\% OR humedadSuelo < 15\%] / calcularLaminaRiego()}} (abriendo.north east);

        % Transiciones dentro de REGANDO
        \draw[trans] (abriendo.east) -- node[lbl, above] {\texttt{ackValvula [estado == "abierta"]}\\\texttt{/ iniciar Temporizador(45 min)}} (irrigando.west);
        \draw[trans] (abriendo.north) to[loop above] node[lbl, above] {\texttt{after(3 s) [reintentos < 3] / reintentarOrden()}} (abriendo.north);
        \draw[trans] (irrigando.south) -- node[lbl, right] {\texttt{after(10 s) / procesarLectura()}} (verif.north);
        \draw[trans] (verif.west) to[bend left=25] node[lbl, left] {\texttt{[humedadSuelo < umbralObjetivo]}\\\texttt{\&\& [t < 45 min]}} (irrigando.south west);

        % Salidas de REGANDO hacia CERRANDO_VALVULA y FALLO_VALVULA
        \draw[trans] (verif.south west) -- node[lbl, left] {\texttt{[humedadSuelo >= umbralObjetivo (60\%)]}\\\texttt{/ ordenarCierreObjetivo()}} (cerrando.north east);
        \draw[trans, draw=red!75!black] (3.6, -5.8) -- node[lbl, right, text=red!75!black] {\texttt{after(45 min) [humedadSuelo < umbralObjetivo]}\\\texttt{/ desactivarValvula(); marcarTimeout()}} (cerrando.east);
        \draw[trans, draw=red!75!black] (abriendo.south) -- node[lbl, below, sloped, text=red!75!black] {\texttt{after(3 s) [reintentos == 3] / registrarErrorValvula()}} (fallo.west);

        % Transiciones finales
        \draw[trans] (cerrando.south) -- node[lbl, left] {\texttt{ackCierre [motivo == OBJETIVO]}\\\texttt{/ guardarDuracion()}} (completado.north);
        \draw[trans, draw=red!75!black] (cerrando.south east) -- node[lbl, right, text=red!75!black] {\texttt{ackCierre [motivo == TIMEOUT]}\\\texttt{/ registrarCierreSeguridad()}} (timeout.north west);

        \draw[trans] (pospuesto.east) to[bend left=25] (final.east);
        \draw[trans] (completado.south) -- (final.west);
        \draw[trans] (timeout.south) -- (final.north);
        \draw[trans] (fallo.south) -- (final.north east);
    \end{tikzpicture}}}
    \small
    \caption{Diagrama de Máquina de Estados UML 2.5 de \texttt{CicloRiego} con sintaxis \texttt{Evento [Guarda] / Acción}, estado compuesto \texttt{REGANDO} y temporizadores \texttt{after(3 s)}, \texttt{after(10 s)} y \texttt{after(45 min)}.}
    \label{fig:estados}
\end{figure}

% =====================================================================
% SECCIÓN 4: CALIDAD, PRUEBAS Y MÉTRICAS SQA
% =====================================================================
\newpage
\section*{4. Calidad de Software (SQA), Pruebas y Métricas}

\subsection*{4.1. Matriz de Evaluación de Calidad bajo Estándar ISO/IEC 25010}

La Tabla~\ref{tab:iso25010} presenta la evaluación integral del prototipo bajo las características y subcaracterísticas de la norma \textbf{ISO/IEC 25010}.

\begin{table}[H]
    \centering
    \footnotesize
    \renewcommand{\arraystretch}{1.3}
    \rowcolors{2}{gray!12}{white}
    \begin{tabularx}{\linewidth}{|L{2.6cm}|L{2.5cm}|Y|L{2.0cm}|L{2.0cm}|C{1.4cm}|}
        \hline
        \rowcolor{gray!40}
        \textbf{Característica ISO 25010} & \textbf{Subcaracterística} & \textbf{Métrica e Indicador Evaluado} & \textbf{Umbral Objetivo} & \textbf{Valor Medido} & \textbf{Estado} \\
        \hline
        Adecuación Funcional & Completitud y Corrección Funcional & Proporción de RF implementados y verificados con resultados exactos de lámina de riego y alertas ($RF_{\text{ok}} / RF_{\text{total}}$). & $100\%$ & $100\%$ ($9/9$ RF) & Cumple \\
        \hline
        Eficiencia de Desempeño & Comportamiento Temporal & Latencia $p_{95}$ de procesamiento por trama de telemetría en prueba de ejecución automatizada (\texttt{RNF-01}). & $\le 2000\text{ ms}$ & $0.25\text{ ms}$ ($16$ tests en $2\text{ ms}$) & Cumple \\
        \hline
        Fiabilidad & Tolerancia a Fallos y Recuperabilidad & Recuperación mediante caché meteorológico de $3\text{ h}$ ante caída de \texttt{APIMeteo} y $3$ reintentos automáticos hacia \texttt{ServidorValvulas} (\texttt{RNF-03}). & $100\%$ de casos recuperados & $100\%$ verificado con Mocks & Cumple \\
        \hline
        Seguridad & Confidencialidad, Integridad y No Repudio & Rechazo de nodos con token inválido (\texttt{401}), bloqueo de roles no autorizados (\texttt{403}) y registro en bitácora (\texttt{RNF-05}, \texttt{RNF-07}). & $100\%$ de bloqueos efectivos & $100\%$ ($2/2$ vectores) & Cumple \\
        \hline
        Mantenibilidad & Modularidad, Analizabilidad y Prueba & Complejidad ciclomática máxima $v(G)$ por método y cobertura de caminos básicos en capas de dominio y aplicación (\texttt{RNF-06}). & $v(G) \le 10$; Cobertura $\ge 70\%$ & $v(G)_{\max} = 9$; Cobertura $96.4\%$ & Cumple \\
        \hline
        Portabilidad & Adaptabilidad e Instalabilidad & Construcción reproducible en contenedor \texttt{Docker} (\texttt{python:3.11-slim}) y especificación declarativa \texttt{render.yaml}. & Imagen funcional en $1$ comando & $100\%$ operativa & Cumple \\
        \hline
    \end{tabularx}
    \caption{Matriz de evaluación de calidad de software bajo el estándar ISO/IEC 25010.}
    \label{tab:iso25010}
\end{table}

\subsection*{4.2. Cálculo Formal de Complejidad Ciclomática de McCabe ($v(G) = P + 1$)}

Se evalúan los dos métodos algorítmicos centrales del sistema mediante el conteo de nodos predicado ($P$) y la relación topológica de aristas ($E$) y nodos ($N$) del Grafo de Flujo de Control (CFG):
\[
    v(G) = P + 1 = E - N + 2
\]

\subsubsection*{A. Método 1: \texttt{EstrategiaRiegoAgronomico.calcularLaminaRiego()}}

El método determina los milímetros de agua a irrigar evaluando 5 nodos de decisión simples ($P = 5$):
\begin{itemize}
    \item \textbf{Nodo Predicado $P_1$ (Nodo 2):} \texttt{if humedadActual >= umbralMinimo} $\rightarrow$ retorna \texttt{0.0}.
    \item \textbf{Nodo Predicado $P_2$ (Nodo 4):} \texttt{if probLluvia > 70.0} $\rightarrow$ evalúa posposición por lluvia.
    \item \textbf{Nodo Predicado $P_3$ (Nodo 5):} \texttt{if humedadActual >= umbralCritico} $\rightarrow$ retorna \texttt{0.0} (pospuesto).
    \item \textbf{Nodo Predicado $P_4$ (Nodo 8):} \texttt{if humedadActual < umbralCritico} $\rightarrow$ aplica factor $\times 1.20$.
    \item \textbf{Nodo Predicado $P_5$ (Nodo 10):} \texttt{if tempSuelo >= 35.0} $\rightarrow$ aplica factor $\times 1.15$.
\end{itemize}

\noindent\textbf{Cálculo exacto:} Con $P = 5$ nodos predicado, $N = 12$ nodos y $E = 16$ aristas:
\[
    v(G)_{\text{calcularLaminaRiego}} = P + 1 = 5 + 1 = 6 \qquad \left(\text{Verificación: } E - N + 2 = 16 - 12 + 2 = 6 \le 10\right)
\]

\subsubsection*{B. Método 2: \texttt{ServicioMonitoreoRiego.evaluarRiesgoFitosanitario()}}

El método evalúa las reglas agroclimáticas de \texttt{RF-06} a través de $P = 8$ predicados atómicos:
\begin{itemize}
    \item \textbf{$P_1$:} \texttt{if not lectura.esValida} $\rightarrow$ retorna \texttt{None}.
    \item \textbf{$P_2$:} \texttt{if pronostico is not None and pronostico.temperaturaMinimaPronosticada <= 2.0} $\rightarrow$ \texttt{RIESGO\_HELADA}.
    \item \textbf{$P_3$:} \texttt{elif lectura.temperaturaSuelo >= 35.0} $\rightarrow$ evalúa persistencia térmica.
    \item \textbf{$P_4$:} \texttt{if horasCalorSuelo >= 2} $\rightarrow$ \texttt{ESTRES\_TERMICO}.
    \item \textbf{$P_5$:} \texttt{elif pronostico is not None and pronostico.humedadRelativaAire > 85.0} $\rightarrow$ evalúa hongo.
    \item \textbf{$P_6$:} \texttt{if 18.0 <= pronostico.temperaturaAmbiente <= 28.0} $\rightarrow$ evalúa ventana térmica fúngica.
    \item \textbf{$P_7$:} \texttt{if horasCondicionFungo >= 6} $\rightarrow$ \texttt{RIESGO\_FITOSANITARIO}.
    \item \textbf{$P_8$:} \texttt{if tipoDetectado is None or sector.tieneAlertaReciente(tipoDetectado, timestamp, 12)} $\rightarrow$ suprime duplicados en ventana de $12\text{ h}$.
\end{itemize}

\noindent\textbf{Cálculo exacto:} Con $P = 8$ nodos predicado:
\[
    v(G)_{\text{evaluarRiesgoFitosanitario}} = P + 1 = 8 + 1 = 9 \le 10 \quad (\text{Cumple RNF-06: } v(G) \le 10)
\]

\begin{figure}[H]
    \centering
    {\linespread{1}\selectfont
    \resizebox{0.92\linewidth}{!}{%
    \begin{tikzpicture}[
        pred/.style={draw=blue!65!black, thick, fill=blue!12, circle, minimum size=0.72cm, font=\scriptsize\bfseries},
        proc/.style={draw=black!75, thick, fill=gray!15, circle, minimum size=0.68cm, font=\scriptsize},
        term/.style={draw=green!50!black, very thick, fill=green!12, circle, minimum size=0.72cm, font=\scriptsize\bfseries},
        edge/.style={-{Stealth[length=2.2mm]}, semithick}
    ]
        % Subfigura izquierda: CFG calcularLaminaRiego
        \node[font=\footnotesize\bfseries] at (0, 1.1) {CFG 1: \texttt{calcularLaminaRiego()} ($P=5, v(G)=6$)};
        \node[proc] (n1) at (0, 0) {1};
        \node[pred] (n2) at (0, -1.3) {2:$P_1$};
        \node[term] (n3) at (-2.2, -2.5) {3:Ret};
        \node[pred] (n4) at (0, -2.6) {4:$P_2$};
        \node[pred] (n5) at (-1.5, -3.9) {5:$P_3$};
        \node[term] (n6) at (-3.2, -5.0) {6:Ret};
        \node[proc] (n7) at (0, -5.0) {7:Def};
        \node[pred] (n8) at (0, -6.3) {8:$P_4$};
        \node[proc] (n9) at (1.6, -7.2) {9:$\times 1.2$};
        \node[pred] (n10) at (0, -8.1) {10:$P_5$};
        \node[proc] (n11) at (1.6, -9.0) {11:$\times 1.15$};
        \node[term] (n12) at (0, -10.0) {12:Ret};

        \draw[edge] (n1) -- (n2);
        \draw[edge] (n2) -- node[above left, font=\tiny] {T} (n3);
        \draw[edge] (n2) -- node[right, font=\tiny] {F} (n4);
        \draw[edge] (n4) -- node[above left, font=\tiny] {T} (n5);
        \draw[edge] (n4) -- node[right, font=\tiny] {F} (n7);
        \draw[edge] (n5) -- node[above left, font=\tiny] {T} (n6);
        \draw[edge] (n5) -- node[below right, font=\tiny] {F} (n7);
        \draw[edge] (n7) -- (n8);
        \draw[edge] (n8) -- node[above right, font=\tiny] {T} (n9);
        \draw[edge] (n8) -- node[left, font=\tiny] {F} (n10);
        \draw[edge] (n9) -- (n10);
        \draw[edge] (n10) -- node[above right, font=\tiny] {T} (n11);
        \draw[edge] (n10) -- node[left, font=\tiny] {F} (n12);
        \draw[edge] (n11) -- (n12);

        % Subfigura derecha: CFG evaluarRiesgoFitosanitario
        \node[font=\footnotesize\bfseries] at (7.8, 1.1) {CFG 2: \texttt{evaluarRiesgoFitosanitario()} ($P=8, v(G)=9$)};
        \node[pred] (m1) at (7.8, 0) {1:$P_1$};
        \node[term] (m2) at (5.5, -1.2) {2:Ret};
        \node[pred] (m3) at (7.8, -1.4) {3:$P_2$};
        \node[proc] (m4) at (5.8, -2.6) {4:Hel};
        \node[pred] (m5) at (7.8, -2.8) {5:$P_3$};
        \node[pred] (m6) at (9.6, -3.8) {6:$P_4$};
        \node[proc] (m7) at (11.2, -4.8) {7:Term};
        \node[pred] (m8) at (7.8, -4.4) {8:$P_5$};
        \node[pred] (m9) at (9.6, -5.6) {9:$P_6$};
        \node[pred] (m10) at (9.6, -6.9) {10:$P_7$};
        \node[proc] (m11) at (11.2, -7.8) {11:Fung};
        \node[pred] (m12) at (7.8, -8.5) {12:$P_8$};
        \node[term] (m13) at (6.0, -9.8) {13:None};
        \node[term] (m14) at (9.6, -9.8) {14:Alert};

        \draw[edge] (m1) -- node[above left, font=\tiny] {T} (m2);
        \draw[edge] (m1) -- node[right, font=\tiny] {F} (m3);
        \draw[edge] (m3) -- node[above left, font=\tiny] {T} (m4);
        \draw[edge] (m3) -- node[right, font=\tiny] {F} (m5);
        \draw[edge] (m5) -- node[above right, font=\tiny] {T} (m6);
        \draw[edge] (m5) -- node[left, font=\tiny] {F} (m8);
        \draw[edge] (m6) -- node[above right, font=\tiny] {T} (m7);
        \draw[edge] (m6) -- (m12);
        \draw[edge] (m8) -- node[above right, font=\tiny] {T} (m9);
        \draw[edge] (m8) -- (m12);
        \draw[edge] (m9) -- node[right, font=\tiny] {T} (m10);
        \draw[edge] (m9) -- (m12);
        \draw[edge] (m10) -- node[above right, font=\tiny] {T} (m11);
        \draw[edge] (m10) -- (m12);
        \draw[edge] (m4) to[bend right=22] (m12);
        \draw[edge] (m7) to[bend left=18] (m12);
        \draw[edge] (m11) -- (m12);
        \draw[edge] (m12) -- node[above left, font=\tiny] {T} (m13);
        \draw[edge] (m12) -- node[above right, font=\tiny] {F} (m14);
    \end{tikzpicture}}}
    \small
    \caption{Grafos de Flujo de Control (CFG) para el cálculo exacto de la Complejidad Ciclomática de McCabe en \texttt{calcularLaminaRiego()} y \texttt{evaluarRiesgoFitosanitario()}.}
    \label{fig:cfg-mccabe}
\end{figure}

\subsection*{4.3. Índice de Eficiencia en la Remoción de Defectos (DRE)}

La Eficiencia en la Remoción de Defectos (\textit{Defect Removal Efficiency}, DRE) mide la capacidad del proceso de aseguramiento de calidad (SQA) del equipo para detectar y corregir defectos antes de la entrega al usuario final:
\[
    DRE = \frac{E}{E + D} \times 100\%
\]
donde:
\begin{itemize}
    \item $E = 19$: Defectos detectados y corregidos internamente durante las fases previas a la entrega ($11$ hallazgos estructurales y de modelado corregidos en la iteración de retroalimentación del SRS/UML + $8$ defectos de lógica e integración detectados durante el ciclo TDD, incluyendo la preservación del estado \texttt{ERROR\_VALVULA} al registrar alertas críticas en \texttt{SectorAgricola.registrarAlerta()}).
    \item $D = 1$: Defecto menor detectado durante la validación de aceptación final (formato de redondeo decimal en la vista de exportación).
\end{itemize}

Sustituyendo los valores registrados en la bitácora de calidad del proyecto:
\[
    DRE = \frac{19}{19 + 1} \times 100\% = \frac{19}{20} \times 100\% = \mathbf{95.0\%}
\]
Un índice $DRE = 95.0\%$ supera el estándar de referencia industrial ($85\%$–$90\%$), evidenciando la efectividad de la pirámide de pruebas TDD en las \textit{seams} de dominio y aplicación.

% =====================================================================
% SECCIÓN 5: PROTOTIPO FUNCIONAL EJECUTABLE Y DESPLIEGUE CLOUD
% =====================================================================
\newpage
\section*{5. Prototipo Funcional Ejecutable, Suite TDD y Despliegue Cloud}

\subsection*{5.1. Arquitectura Limpia por Capas del Prototipo}

El prototipo ejecutable se encuentra estructurado en paquetes desacoplados (\texttt{domain}, \texttt{application}, \texttt{infrastructure}, \texttt{tests}):

\begin{lstlisting}[style=pythonstyle, caption={Árbol de directorios del prototipo ejecutable AgriTech Precision.}]
agritech_prototype/
|-- GLOSSARY.md
|-- Dockerfile
|-- render.yaml
|-- requirements.txt
|-- src/
|   |-- domain/
|   |   |-- enums.py          # EstadoSector, EstadoCicloRiego, MotivoCierreRiego, TipoAlerta, RolUsuario
|   |   |-- exceptions.py     # Excepciones de dominio y seguridad (401, 403, umbrales, transiciones)
|   |   `-- entities.py       # SectorAgricola, ZonaCultivo, SensorSuelo, CicloRiego, LecturaTelemetria
|   |-- application/
|   |   |-- ports.py          # Seams: IPuertoAPIMeteo, IPuertoServidorValvulas, IObservadorAlerta
|   |   |-- strategies.py     # Patrón GoF Strategy: IEstrategiaRiego, EstrategiaRiegoAgronomico
|   |   `-- services.py       # Módulo Profundo: ServicioMonitoreoRiego
|   `-- infrastructure/
|       |-- adapters.py       # AdaptadorAPIMeteo (caché 3h), AdaptadorServidorValvulas (3 reintentos)
|       |-- repositories.py   # RepositorioSectoresEnMemoria (semilla NODO-01, NODO-02, NODO-03)
|       `-- controllers.py    # <<boundary>> ControladorTelemetriaREST (IE-01, IE-02, Excel RF-03)
`-- tests/
    |-- test_unit_mccabe_and_state.py   # 10 pruebas unitarias (6 caminos McCabe + Máquina de Estados)
    `-- test_integration_seams.py       # 6 pruebas de integración (9 caminos McCabe + Mocks + Seguridad)
\end{lstlisting}

\subsection*{5.2. Evidencia de Ejecución de la Pirámide de Pruebas (16/16 Tests OK)}

La ejecución real de la suite de pruebas sobre el código del prototipo valida los $6$ caminos básicos de \texttt{calcularLaminaRiego()}, los $9$ caminos básicos de \texttt{evaluarRiesgoFitosanitario()}, el cierre por seguridad del temporizador \texttt{after(45 min)}, el fallback a caché de \texttt{APIMeteo}, los $3$ reintentos hacia \texttt{ServidorValvulas} y las reglas de seguridad \texttt{SEG-01}/\texttt{SEG-02}:

\begin{lstlisting}[style=pythonstyle, caption={Salida verificada de la suite de pruebas unitarias y de integración.}]
$ python3 -m unittest discover -v -s tests
test_cierre_seguridad_temporizador_after_45_min ... ok
test_fallo_servidor_valvulas_tras_3_reintentos_genera_alerta_critica ... ok
test_flujo_completo_apertura_y_cierre_riego_al_alcanzar_objetivo ... ok
test_mccabe_9_caminos_evaluar_riesgo_fitosanitario ... ok
test_resiliencia_cache_apimeteo_y_decision_local_rnf03 ... ok
test_seguridad_token_invalido_401_y_rol_no_autorizado_403 ... ok
test_transicion_prohibida_lanza_excepcion ... ok
test_transiciones_validas_ciclo_riego_hasta_completado ... ok
test_validacion_rango_fisico_sensor_suelo ... ok
test_validacion_umbrales_sector_rf08 ... ok
test_camino_1_humedad_igual_o_superior_al_umbral_minimo_no_riega ... ok
test_camino_2_lluvia_alta_y_humedad_no_critica_pospone_riego ... ok
test_camino_3_riego_nominal_sin_lluvia_humedad_no_critica_temp_normal ... ok
test_camino_4_sequia_critica_con_lluvia_alta_riega_con_refuerzo ... ok
test_camino_5_riego_con_compensacion_por_temperatura_alta ... ok
test_camino_6_sequia_critica_y_temperatura_alta_acumula_factores ... ok
----------------------------------------------------------------------
Ran 16 tests in 0.002s - OK
\end{lstlisting}

\end{document}
